%% Supplementary Information, typeset with the Springer Nature sn-jnl class.
\documentclass[pdflatex,sn-nature]{sn-jnl}

\usepackage{graphicx}\usepackage{multirow}\usepackage{amsmath,amssymb,amsfonts}
\usepackage{amsthm}\usepackage{array}\usepackage{booktabs}\usepackage{xcolor}

\newcommand{\X}{\mathcal{X}}
\newcommand{\C}{\mathcal{C}}
\newcommand{\F}{\mathcal{F}}
\newcommand{\PS}{\mathcal{P}^{\star}}
\newcommand{\LB}{\mathit{LB}}
\theoremstyle{thmstyleone}
\newtheorem{theorem}{Supplementary Theorem}
\newtheorem{proposition}[theorem]{Supplementary Proposition}
\theoremstyle{thmstyletwo}
\newtheorem{remark}{Remark}
\theoremstyle{thmstylethree}
\newtheorem{definition}{Supplementary Definition}

\renewcommand{\thetable}{S\arabic{table}}
\renewcommand{\thefigure}{S\arabic{figure}}
\renewcommand{\theequation}{S\arabic{equation}}

\raggedbottom

\begin{document}

\title[Supplementary Information]{Supplementary Information:\\
Exact architecture exploration for heterogeneous edge and cloud cyber-physical systems}

\author*[1]{\fnm{Nikita} \sur{Tarasov}}\email{dev.nikita@outlook.com}
\author[1]{\fnm{Roman} \sur{Zinko}}\email{roman.v.zinko@lpnu.ua}
\affil*[1]{\orgdiv{Institute of Computer Science and Information Technologies},
\orgname{Lviv Polytechnic National University}, \orgaddress{\city{Lviv},
\postcode{79013}, \country{Ukraine}}}

\abstract{Supplementary Notes 1--2 and Supplementary Tables S1--S13 accompanying the
main manuscript. All tabulated values are generated from the same result files as the
main-text summaries.}

\maketitle

\section{Supplementary Note 1: correctness of the pruning rules}\label{snote1}

\subsection{Setting}
The design space is the Cartesian product $\X=D_1\times\dots\times D_k$ traversed under a
fixed variable order $v_1,\dots,v_k$. A \emph{partial design} of depth $d$ is a prefix of
assignments $p=(a_1,\dots,a_d)$, $a_i\in D_i$, and its completions are
$\C(p)=\{x\in\X: x_i=a_i,\ i\le d\}$.

Let $\sigma:\X\to\{0,1\}$ be the well-formedness predicate, $\X_\sigma$ the well-formed
space, $\F$ the feasible set and $\F_\sigma=\F\cap\X_\sigma$. Objectives
$F(x)=(f_1(x),\dots,f_n(x))$ are minimised, $u\preceq v$ denotes $u_i\le v_i$ for all $i$
with at least one strict inequality, and
\begin{equation}
\PS=\{F(x): x\in\F_\sigma,\ \nexists\,y\in\F_\sigma:\ F(y)\preceq F(x)\}.
\end{equation}
A bound $\LB_i$ is \emph{admissible} for $f_i$ if $\LB_i(p)\le f_i(x)$ for every partial
design $p$ and every $x\in\C(p)$.

\subsection{Construction of the bounds}

\begin{theorem}[Term-wise relaxation]\label{sthm:relax}
Let $f_i(x)=\sum_{t\in T_i}\varphi_t(x)$ with $\varphi_t\ge0$ and
$\underline{\varphi}_t(p)=\min_{x\in\C(p)}\varphi_t(x)$. Then
$\LB_i(p)=\sum_{t\in T_i}\underline{\varphi}_t(p)$ is admissible.
\end{theorem}
\begin{proof}
For every $x\in\C(p)$ and every $t\in T_i$ we have
$\underline{\varphi}_t(p)\le\varphi_t(x)$, because the minimum is taken over a set that
contains $x$. Summing over $t\in T_i$ gives $\LB_i(p)\le\sum_t\varphi_t(x)=f_i(x)$.
\end{proof}

The minima are taken independently per term, so the combination of values realising them
need not correspond to any single completion; the bound is optimistic but valid. Because the
decision domains are finite, the minimum of each term over the reachable domain is
precomputed exactly and read from a table in constant time; for ordered scalar domains
the minimum is attained at a domain endpoint, which is how the table is built.

\begin{theorem}[Exactness on determined terms]\label{sthm:exact}
Let $\mathrm{supp}(\varphi_t)$ be the set of variables $\varphi_t$ depends on. If
$\mathrm{supp}(\varphi_t)\subseteq\{v_1,\dots,v_d\}$, then $\varphi_t$ is constant on
$\C(p)$, and substituting its exact value for $\underline{\varphi}_t(p)$ preserves
admissibility while tightening the bound.
\end{theorem}
\begin{proof}
Constancy follows from the definition of $\C(p)$, and a constant equals its own minimum
over $\C(p)$, so Theorem~\ref{sthm:relax} applies with the substituted value.
\end{proof}

\begin{remark}[Role of the variable order]
The order places the determinants of tier utilisation --- node counts, node classes,
replication factor and placement policy --- before the remaining variables. Once they are
assigned, the waiting term $\frac{w_i}{C_t}\cdot\frac{\rho_t}{2(1-\rho_t)}$ is constant
on $\C(p)$ and enters the bound exactly by Theorem~\ref{sthm:exact} instead of being
relaxed to its trivial bound $0$. This is what makes a contention-dependent objective
usable for safe pruning without assignment monotonicity. If the determinants are fixed
and $\rho_t\ge1$, no completion has finite latency and $\LB_L(p)=\infty$, which is
admissible.
\end{remark}

\subsection{The three pruning rules}

\begin{proposition}[Safe feasibility pruning]\label{sprop:feas}
If $\LB_{g_j}$ is admissible for the constrained quantity $g_j$ and
$\LB_{g_j}(p)>g_j^{\max}$, then $\C(p)\cap\F=\varnothing$.
\end{proposition}
\begin{proof}
For any $x\in\C(p)$, admissibility gives $g_j(x)\ge\LB_{g_j}(p)>g_j^{\max}$, so $x$
violates the constraint and $x\notin\F$.
\end{proof}

\begin{proposition}[Safe dominance pruning]\label{sprop:dom}
Let $y\in\F_\sigma$ be a fully evaluated feasible architecture with
$F(y)\preceq\LB(p)$. Then $F(x)\notin\PS$ for every $x\in\C(p)$.
\end{proposition}
\begin{proof}
Let $x\in\C(p)$. Admissibility gives $f_i(x)\ge\LB_i(p)\ge F_i(y)$ for all $i$, and by
assumption there is a $j$ with $F_j(y)<\LB_j(p)\le f_j(x)$. Hence $y$ strictly dominates
$x$, so $F(x)$ is not on the Pareto front.
\end{proof}

\begin{proposition}[Closure of the well-formedness predicate]\label{sprop:struct}
Each invariant of $\sigma$ depends only on assigned variables and, under the chosen
order, all variables of an invariant are assigned before it is first checked. Since
assignments are never revised, $\neg\sigma(p)$ implies $\sigma(x)=0$ for all
$x\in\C(p)$, that is $\C(p)\cap\X_\sigma=\varnothing$.
\end{proposition}

\begin{remark}
The structural invariants define what counts as an architecture rather than following
from the constraints of the problem: the evaluator can formally compute objectives for,
say, a replicated deployment under a policy that places the ingest stage off the edge
tier, although such a configuration is meaningless. Structural pruning is therefore part
of the definition of the search space, and the exhaustive baseline applies the same
predicate.
\end{remark}

\subsection{Front preservation}

\begin{theorem}[Pareto-front preservation]\label{sthm:front}
If all bounds are admissible, pruning is performed only by
Propositions~\ref{sprop:feas}--\ref{sprop:struct}, and the archive contains only feasible
fully evaluated architectures, then the set of non-dominated vectors returned by the
exploration equals $\PS$.
\end{theorem}
\begin{proof}
$(\subseteq)$ Every returned vector is $F(x)$ for an evaluated $x\in\F_\sigma$ that
survives non-dominated filtering. A vector dominating it would belong to an architecture
that was either evaluated, in which case filtering would have removed the dominated one,
or pruned; pruning by Proposition~\ref{sprop:feas} or \ref{sprop:struct} would contradict
its feasibility and well-formedness, and pruning by Proposition~\ref{sprop:dom} implies a
further feasible evaluated architecture dominating it, so the argument repeats along a
strictly decreasing chain in a finite set.

$(\supseteq)$ Let $u\in\PS$ and $x\in\F_\sigma$ with $F(x)=u$. The traversal is a prefix
tree, so $x$ is reached through the unique chain
$p_0\subset p_1\subset\dots\subset p_k=x$. If some $p_d$ were pruned, then by
Proposition~\ref{sprop:struct} $\sigma(x)=0$ and by Proposition~\ref{sprop:feas}
$x\notin\F$, both contradicting $x\in\F_\sigma$; by Proposition~\ref{sprop:dom} there
would exist a feasible $y$ strictly dominating $x$, contradicting $u\in\PS$. Hence no
$p_d$ is pruned, $x$ is evaluated, and $u$ survives the final filtering.
\end{proof}

\begin{remark}[Duplicates]
The theorem is stated in objective space. If two distinct architectures share an
objective vector, dominance pruning may retain only one of them; the front is unaffected.
All experiments therefore compare sets of objective vectors.
\end{remark}

\begin{remark}[Order independence]
Proposition~\ref{sprop:dom} holds for any archive of evaluated feasible architectures, so
the traversal order affects pruning efficiency but never correctness. This is confirmed
empirically in Supplementary Table~S7, where all order and toggle combinations return the
same front.
\end{remark}

\subsection{Complexity and break-even}
With $C_E$ the cost of one full evaluation, $C_P$ the cost of computing the bounds at one
node, $N$ the number of expanded partial designs, $N_E$ the number of full evaluations
and $|A|$ the archive size,
\begin{equation}
T_{\mathrm{exh}}=O\big(|\X|\,C_E\big),\qquad
T_{\mathrm{HCA}}=O\big(N(C_P+n|A|)+N_E\,C_E\big),
\end{equation}
so pruning is profitable when $C_E>N(C_P+n|A|)/(|\X|-N_E)$. The right-hand side is
measured rather than assumed. In the runs reported here it is negative on every scale and
workload, because the total search overhead $N(C_P+n|A|)$ is already smaller than the
overhead of enumerating the space: there is no positive evaluation cost below which
exhaustive enumeration wins, and the advantage grows from there with $C_E$.

\section{Supplementary Note 2: full evaluation model}\label{snote2}

With $\mathcal{H}(x)$ the hop set implied by the placement policy and the replication
factor, $\Lambda$ the system arrival rate, $T$ the accounting horizon and $\rho_t$ the
tier utilisation defined in Methods:
\begin{align}
E(x) &= \sum_{t\in\{e,g\}} N_t\Big(P^{\mathrm{idle}}_t+\rho_t\big(P^{\mathrm{peak}}_t-
        P^{\mathrm{idle}}_t\big)\Big)T
      + \rho_c P^{\mathrm{peak}}_c T
      + \sum_{(s,\ell)\in\mathcal{H}(x)} s\Lambda e_\ell T, \\
C(x) &= n_e\big(c_{\kappa_e}+c_\ell\big)+n_g\big(c_{\kappa_g}+c_\ell\big)+c_{\kappa_c}, \\
D(x) &= \Lambda \sum_{(s,\ell)\in\mathcal{H}(x)} s .
\end{align}
An architecture is feasible when $\rho_t\le0.85$ on every used tier, the resident memory
of the stages mapped to a tier fits the memory of its node class, the aggregate
access-link traffic does not exceed the total bandwidth of the attached nodes, and the
latency, cost and energy budgets hold. Energy is reported in kWh over $T=1$\,h. Cost
coefficients are synthetic normalised values encoding relative price tiers; no monetary
interpretation is claimed, and neither the energy nor the cost model is validated against
measurements.

\section{Supplementary Tables}\label{stables}

\begin{table}[!htbp]\centering
\caption{Full budget-matched comparison on S2. Medians over 20 seeds.}\label{tabS1}
{\let\scriptsize\tiny\setlength{\tabcolsep}{2pt}\scriptsize
\begin{tabular}{lllrrrrr}
\toprule
scale & workload & method & evals & recall & HV ratio & IGD$^{+}$ & time [s] \\
\midrule
\multirow{10}{*}{S2} & \multirow{10}{*}{telemetry} & exhaustive & 3\,888 & 1.000 & 1.000 & 0.0000 & 0.045 \\
 & & random ($B$=100) & 100 & 0.037 & 0.692 & 0.1496 & 0.001 \\
 & & random ($B$=250) & 250 & 0.037 & 0.794 & 0.1018 & 0.003 \\
 & & random ($B$=500) & 500 & 0.074 & 0.880 & 0.0713 & 0.007 \\
 & & random ($B$=1000) & 1000 & 0.185 & 0.923 & 0.0451 & 0.014 \\
 & & NSGA-II ($B$=100) & 100 & 0.037 & 0.703 & 0.1537 & 0.011 \\
 & & NSGA-II ($B$=250) & 250 & 0.259 & 0.886 & 0.0447 & 0.029 \\
 & & NSGA-II ($B$=500) & 500 & 0.759 & 0.995 & 0.0063 & 0.072 \\
 & & NSGA-II ($B$=1000) & 1000 & 0.944 & 0.999 & 0.0003 & 0.148 \\
 & & \textbf{HCA-DSE} & 46 & 1.000 & 1.000 & 0.0000 & 0.013 \\
\midrule
\multirow{10}{*}{S2} & \multirow{10}{*}{control} & exhaustive & 3\,888 & 1.000 & 1.000 & 0.0000 & 0.049 \\
 & & random ($B$=100) & 100 & 0.040 & 0.691 & 0.1460 & 0.001 \\
 & & random ($B$=250) & 250 & 0.040 & 0.792 & 0.0940 & 0.003 \\
 & & random ($B$=500) & 500 & 0.080 & 0.879 & 0.0566 & 0.006 \\
 & & random ($B$=1000) & 1000 & 0.160 & 0.923 & 0.0379 & 0.014 \\
 & & NSGA-II ($B$=100) & 100 & 0.040 & 0.627 & 0.1592 & 0.011 \\
 & & NSGA-II ($B$=250) & 250 & 0.300 & 0.910 & 0.0348 & 0.030 \\
 & & NSGA-II ($B$=500) & 500 & 0.800 & 0.991 & 0.0041 & 0.072 \\
 & & NSGA-II ($B$=1000) & 1000 & 0.900 & 0.999 & 0.0007 & 0.161 \\
 & & \textbf{HCA-DSE} & 39 & 1.000 & 1.000 & 0.0000 & 0.014 \\
\midrule
\multirow{10}{*}{S2} & \multirow{10}{*}{sensing} & exhaustive & 3\,888 & 1.000 & 1.000 & 0.0000 & 0.029 \\
 & & random ($B$=100) & 100 & 0.000 & 0.469 & 0.2629 & 0.001 \\
 & & random ($B$=250) & 250 & 0.062 & 0.686 & 0.1521 & 0.002 \\
 & & random ($B$=500) & 500 & 0.062 & 0.804 & 0.0808 & 0.005 \\
 & & random ($B$=1000) & 1000 & 0.188 & 0.866 & 0.0622 & 0.010 \\
 & & NSGA-II ($B$=100) & 100 & 0.000 & 0.404 & 0.3482 & 0.010 \\
 & & NSGA-II ($B$=250) & 250 & 0.312 & 0.793 & 0.0782 & 0.031 \\
 & & NSGA-II ($B$=500) & 500 & 0.719 & 0.965 & 0.0224 & 0.077 \\
 & & NSGA-II ($B$=1000) & 1000 & 0.938 & 1.000 & 0.0009 & 0.186 \\
 & & \textbf{HCA-DSE} & 26 & 1.000 & 1.000 & 0.0000 & 0.010 \\
\midrule
\bottomrule
\end{tabular}}
\end{table}

\begin{table}[!htbp]\centering
\caption{Full budget-matched comparison on S3. Medians over 20 seeds.}\label{tabS2}
{\let\scriptsize\tiny\setlength{\tabcolsep}{2pt}\scriptsize
\begin{tabular}{lllrrrrr}
\toprule
scale & workload & method & evals & recall & HV ratio & IGD$^{+}$ & time [s] \\
\midrule
\multirow{10}{*}{S3} & \multirow{10}{*}{telemetry} & exhaustive & 34\,560 & 1.000 & 1.000 & 0.0000 & 0.320 \\
 & & random ($B$=100) & 100 & 0.000 & 0.259 & 0.3903 & 0.001 \\
 & & random ($B$=250) & 250 & 0.000 & 0.487 & 0.2574 & 0.002 \\
 & & random ($B$=500) & 500 & 0.000 & 0.563 & 0.1671 & 0.005 \\
 & & random ($B$=1000) & 1000 & 0.016 & 0.699 & 0.1147 & 0.011 \\
 & & NSGA-II ($B$=100) & 100 & 0.000 & 0.151 & 0.4919 & 0.009 \\
 & & NSGA-II ($B$=250) & 250 & 0.000 & 0.687 & 0.1732 & 0.026 \\
 & & NSGA-II ($B$=500) & 500 & 0.355 & 0.915 & 0.0558 & 0.061 \\
 & & NSGA-II ($B$=1000) & 1000 & 0.645 & 0.987 & 0.0415 & 0.146 \\
 & & \textbf{HCA-DSE} & 52 & 1.000 & 1.000 & 0.0000 & 0.055 \\
\midrule
\multirow{10}{*}{S3} & \multirow{10}{*}{control} & exhaustive & 34\,560 & 1.000 & 1.000 & 0.0000 & 0.276 \\
 & & random ($B$=100) & 100 & 0.000 & 0.262 & 0.4192 & 0.001 \\
 & & random ($B$=250) & 250 & 0.000 & 0.492 & 0.2537 & 0.002 \\
 & & random ($B$=500) & 500 & 0.000 & 0.564 & 0.1746 & 0.005 \\
 & & random ($B$=1000) & 1000 & 0.018 & 0.691 & 0.1248 & 0.010 \\
 & & NSGA-II ($B$=100) & 100 & 0.000 & 0.286 & 0.4405 & 0.010 \\
 & & NSGA-II ($B$=250) & 250 & 0.037 & 0.748 & 0.1009 & 0.025 \\
 & & NSGA-II ($B$=500) & 500 & 0.278 & 0.933 & 0.0387 & 0.063 \\
 & & NSGA-II ($B$=1000) & 1000 & 0.648 & 0.991 & 0.0244 & 0.147 \\
 & & \textbf{HCA-DSE} & 43 & 1.000 & 1.000 & 0.0000 & 0.053 \\
\midrule
\multirow{10}{*}{S3} & \multirow{10}{*}{sensing} & exhaustive & 34\,560 & 1.000 & 1.000 & 0.0000 & 0.282 \\
 & & random ($B$=100) & 100 & 0.000 & 0.037 & 0.6082 & 0.001 \\
 & & random ($B$=250) & 250 & 0.000 & 0.231 & 0.3843 & 0.002 \\
 & & random ($B$=500) & 500 & 0.000 & 0.556 & 0.2543 & 0.005 \\
 & & random ($B$=1000) & 1000 & 0.000 & 0.600 & 0.1747 & 0.010 \\
 & & NSGA-II ($B$=100) & 100 & 0.000 & 0.050 & 0.5191 & 0.011 \\
 & & NSGA-II ($B$=250) & 250 & 0.037 & 0.279 & 0.2833 & 0.025 \\
 & & NSGA-II ($B$=500) & 500 & 0.370 & 0.750 & 0.1996 & 0.064 \\
 & & NSGA-II ($B$=1000) & 1000 & 0.741 & 0.969 & 0.1670 & 0.165 \\
 & & \textbf{HCA-DSE} & 43 & 1.000 & 1.000 & 0.0000 & 0.071 \\
\bottomrule
\end{tabular}}
\end{table}

\begin{table}[!htbp]\centering
\caption{Full S3 statistical comparison: one-sided Mann--Whitney test for larger NSGA-II
hypervolume and Cliff's $\delta$. Unadjusted $p$-values, exploratory.}\label{tabS3}
{\let\scriptsize\tiny\setlength{\tabcolsep}{2pt}\scriptsize
\setlength{\tabcolsep}{3.5pt}
\begin{tabular}{lrrrrrrr}
\toprule
workload & budget & NSGA-II & random & $p$ & Cliff's & NSGA-II & HCA-DSE \\
 & & HV med. & HV med. & & $\delta$ & recall & evals \\
\midrule
telemetry & 100 & 0.1509 & 0.259 & 0.951 & -0.302 & 0.0 & 52 \\
telemetry & 250 & 0.6868 & 0.4871 & 0.021 & 0.38 & 0.0 & 52 \\
telemetry & 500 & 0.9154 & 0.5633 & $<$0.001 & 0.7 & 0.355 & 52 \\
telemetry & 1000 & 0.9874 & 0.6993 & $<$0.001 & 0.7 & 0.645 & 52 \\
control & 100 & 0.2859 & 0.2616 & 0.484 & 0.01 & 0.0 & 43 \\
control & 250 & 0.7481 & 0.4916 & $<$0.001 & 0.62 & 0.037 & 43 \\
control & 500 & 0.9331 & 0.5644 & $<$0.001 & 0.8 & 0.278 & 43 \\
control & 1000 & 0.9906 & 0.6913 & $<$0.001 & 0.8 & 0.648 & 43 \\
sensing & 100 & 0.0497 & 0.0367 & 0.261 & 0.12 & 0.0 & 43 \\
sensing & 250 & 0.2787 & 0.2308 & 0.672 & -0.08 & 0.037 & 43 \\
sensing & 500 & 0.7504 & 0.5561 & 0.168 & 0.18 & 0.37 & 43 \\
sensing & 1000 & 0.9689 & 0.5998 & 0.054 & 0.3 & 0.741 & 43 \\
\bottomrule
\end{tabular}}
\end{table}

The main text compares analytical predictions with the mean latency measured over each
deployment run. Supplementary Table~S4 is the complementary robustness analysis, in which
the same predictions are compared with the measured median ($p_{50}$) of the same runs.

\begin{table}[!htbp]\centering
\caption{Testbed validation against measured $p_{50}$, including utilisation and
reproducibility subsets. Pair agreement is a fraction; undefined values (no eligible
pairs) are shown as --.}\label{tabS4}
{\let\scriptsize\tiny\setlength{\tabcolsep}{2pt}\scriptsize
\setlength{\tabcolsep}{3pt}
\begin{tabular}{lrrrrrrrr}
\toprule
subset & $n$ & MAPE & med.\ AE & Spearman & Kendall & pair & pair & repeat \\
 & & [\%] & [ms] & $\rho$ & $\tau$ & $>10\%$ & $>25\%$ & IQR [ms] \\
\midrule
telemetry (all) & 12 & 19.9 & 7.43 & 0.811 & 0.697 & 0.815 & 0.895 & 0.23 \\
telemetry ($\rho\le0.75$) & 6 & 15.4 & 2.85 & 0.829 & 0.733 & 0.818 & 1.0 & 0.2 \\
control (all) & 12 & 14.4 & 2.49 & 0.865 & 0.748 & 0.896 & 0.909 & 0.11 \\
control ($\rho\le0.75$) & 10 & 9.4 & 2.32 & 0.766 & 0.629 & 0.815 & 0.833 & 0.11 \\
sensing (all) & 12 & 51.1 & 19.15 & 0.935 & 0.84 & 0.935 & 0.913 & 1.0 \\
sensing ($\rho\le0.75$) & 10 & 51.0 & 18.82 & 0.985 & 0.944 & 1.0 & 1.0 & 0.75 \\
sensing (reproducible) & 8 & 39.3 & 18.56 & 0.97 & 0.909 & 1.0 & -- & 0.67 \\
\textbf{all workloads} (all) & 36 & 28.4 & 7.43 & 0.968 & 0.884 & 0.961 & 0.98 & 0.27 \\
\textbf{all workloads} ($\rho\le0.75$) & 26 & 26.8 & 3.96 & 0.945 & 0.855 & 0.959 & 0.982 & 0.22 \\
\textbf{all workloads} (reproducible) & 32 & 22.7 & 4.5 & 0.957 & 0.869 & 0.956 & 0.979 & 0.23 \\
\bottomrule
\end{tabular}}
\end{table}

\begin{table}[!htbp]\centering
\caption{Exact solver baselines against HCA-DSE on S1--S3. Finite-table SMT evaluates
every structurally valid design and then solves; factorised SMT tabulates each model term
over its own decisions only and calls no evaluator. ``Evaluator calls'' counts full
design-point evaluations; ``tabulated entries'' counts table rows built before solving;
total time includes table construction. All rows return the exhaustive front.
Factorised SMT requires the complete analytical objective and constraint model to be
representable by the factor tables; solver times depend on the encoding.}\label{tabS5}
{\let\scriptsize\tiny\setlength{\tabcolsep}{2pt}\scriptsize
\begin{tabular}{lllrrrrr}
\toprule
scale & workload & method & evaluator calls & tabulated entries & solver checks & total [s] & front \\
\midrule
S1 & telemetry & HCA-DSE & 32 & -- & -- & 0.005 & 17 \\
 &  & finite-table SMT & 1080 & 1080 & 44 & 0.341 & 17 \\
 &  & factorised SMT & 0 & 221 & 51 & 1.518 & 17 \\
S1 & control & HCA-DSE & 28 & -- & -- & 0.005 & 18 \\
 &  & finite-table SMT & 1080 & 1080 & 48 & 0.317 & 18 \\
 &  & factorised SMT & 0 & 225 & 52 & 1.684 & 18 \\
S1 & sensing & HCA-DSE & 16 & -- & -- & 0.004 & 11 \\
 &  & finite-table SMT & 1080 & 1080 & 31 & 0.155 & 11 \\
 &  & factorised SMT & 0 & 199 & 37 & 0.766 & 11 \\
S2 & telemetry & HCA-DSE & 46 & -- & -- & 0.014 & 27 \\
 &  & finite-table SMT & 3888 & 3888 & 71 & 3.305 & 27 \\
 &  & factorised SMT & 0 & 529 & 81 & 10.613 & 27 \\
S2 & control & HCA-DSE & 39 & -- & -- & 0.013 & 25 \\
 &  & finite-table SMT & 3888 & 3888 & 61 & 2.180 & 25 \\
 &  & factorised SMT & 0 & 537 & 83 & 12.447 & 25 \\
S2 & sensing & HCA-DSE & 26 & -- & -- & 0.011 & 16 \\
 &  & finite-table SMT & 3888 & 3888 & 48 & 0.937 & 16 \\
 &  & factorised SMT & 0 & 469 & 50 & 7.373 & 16 \\
S3 & telemetry & HCA-DSE & 52 & -- & -- & 0.058 & 31 \\
 &  & factorised SMT & 0 & 2384 & 87 & 80.961 & 31 \\
S3 & control & HCA-DSE & 43 & -- & -- & 0.055 & 27 \\
 &  & factorised SMT & 0 & 2406 & 85 & 94.321 & 27 \\
S3 & sensing & HCA-DSE & 43 & -- & -- & 0.067 & 27 \\
 &  & factorised SMT & 0 & 2188 & 85 & 71.363 & 27 \\
\bottomrule
\end{tabular}}
\end{table}

\begin{table}[!htbp]\centering
\caption{Ablation on S3: expanded partial designs, full evaluations and runtime for each
workload.}\label{tabS6}
{\let\scriptsize\tiny\setlength{\tabcolsep}{2pt}\scriptsize
\setlength{\tabcolsep}{3.5pt}
\begin{tabular}{lrrrrrrrrr}
\toprule
& \multicolumn{3}{c}{telemetry} & \multicolumn{3}{c}{control} & \multicolumn{3}{c}{sensing} \\
\cmidrule(lr){2-4}\cmidrule(lr){5-7}\cmidrule(lr){8-10}
variant & states & evals & [ms] & states & evals & [ms] & states & evals & [ms] \\
\midrule
structural only & 61\,225 & 34\,560 & 369 & 61\,225 & 34\,560 & 284 & 61\,225 & 34\,560 & 284 \\
structural + feasibility bounds & 24\,650 & 12\,700 & 359 & 15\,979 & 7\,025 & 262 & 23\,118 & 10\,956 & 315 \\
structural + dominance bounds & 4\,617 & 323 & 105 & 4\,730 & 405 & 100 & 4\,804 & 412 & 89 \\
\textbf{HCA-DSE (all layers)} & 2\,397 & 52 & 54 & 2\,375 & 43 & 51 & 3\,221 & 43 & 61 \\
\bottomrule
\end{tabular}}
\end{table}

\begin{table}[!htbp]\centering
\caption{Full order and queue-bound ablation on S1--S3. All configurations preserve the
unrounded exhaustive front.}\label{tabS7}
{\let\scriptsize\tiny\setlength{\tabcolsep}{2pt}\scriptsize
\begin{tabular}{llllrrrrrr}
\toprule
scale & workload & order & queue & states & checks & evals & active & pruned & ms \\
\midrule
S1 & telemetry & determinants-early & off & 447 & 746 & 136 & 0 & 0 & 6.7 \\
S1 & telemetry & determinants-early & on & 249 & 516 & 32 & 306 & 46 & 5.2 \\
S1 & telemetry & determinants-late & off & 695 & 1299 & 136 & 0 & 0 & 9.5 \\
S1 & telemetry & determinants-late & on & 589 & 1299 & 30 & 0 & 0 & 10.3 \\
S1 & control & determinants-early & off & 391 & 669 & 113 & 0 & 0 & 5.2 \\
S1 & control & determinants-early & on & 242 & 512 & 28 & 312 & 35 & 4.6 \\
S1 & control & determinants-late & off & 679 & 1317 & 112 & 0 & 0 & 9.6 \\
S1 & control & determinants-late & on & 593 & 1317 & 26 & 0 & 0 & 10.5 \\
S1 & sensing & determinants-early & off & 481 & 810 & 147 & 0 & 0 & 5.8 \\
S1 & sensing & determinants-early & on & 226 & 503 & 16 & 309 & 72 & 4.1 \\
S1 & sensing & determinants-late & off & 497 & 819 & 148 & 0 & 0 & 6.1 \\
S1 & sensing & determinants-late & on & 365 & 819 & 16 & 0 & 0 & 5.8 \\
S2 & telemetry & determinants-early & off & 939 & 1815 & 208 & 0 & 0 & 15.4 \\
S2 & telemetry & determinants-early & on & 604 & 1383 & 46 & 839 & 99 & 13.9 \\
S2 & telemetry & determinants-late & off & 1816 & 3940 & 208 & 0 & 0 & 30.7 \\
S2 & telemetry & determinants-late & on & 1652 & 3940 & 44 & 0 & 0 & 35.8 \\
S2 & control & determinants-early & off & 835 & 1674 & 165 & 0 & 0 & 13.9 \\
S2 & control & determinants-early & on & 587 & 1364 & 39 & 842 & 80 & 12.0 \\
S2 & control & determinants-late & off & 1780 & 3960 & 164 & 0 & 0 & 29.0 \\
S2 & control & determinants-late & on & 1653 & 3960 & 37 & 0 & 0 & 32.9 \\
S2 & sensing & determinants-early & off & 1161 & 2184 & 289 & 0 & 0 & 15.0 \\
S2 & sensing & determinants-early & on & 600 & 1427 & 26 & 879 & 177 & 10.1 \\
S2 & sensing & determinants-late & off & 1268 & 2421 & 290 & 0 & 0 & 16.8 \\
S2 & sensing & determinants-late & on & 1004 & 2421 & 26 & 0 & 0 & 15.8 \\
S3 & telemetry & determinants-early & off & 2871 & 6806 & 251 & 0 & 0 & 53.3 \\
S3 & telemetry & determinants-early & on & 2397 & 6100 & 52 & 3777 & 188 & 54.4 \\
S3 & telemetry & determinants-late & off & 6029 & 14794 & 251 & 0 & 0 & 117.9 \\
S3 & telemetry & determinants-late & on & 5828 & 14794 & 50 & 0 & 0 & 131.9 \\
S3 & control & determinants-early & off & 2755 & 6636 & 206 & 0 & 0 & 51.0 \\
S3 & control & determinants-early & on & 2375 & 6074 & 43 & 3777 & 163 & 56.3 \\
S3 & control & determinants-late & off & 5990 & 14812 & 205 & 0 & 0 & 106.7 \\
S3 & control & determinants-late & on & 5826 & 14812 & 41 & 0 & 0 & 127.5 \\
S3 & sensing & determinants-early & off & 4351 & 9905 & 537 & 0 & 0 & 72.9 \\
S3 & sensing & determinants-early & on & 3221 & 8281 & 43 & 5139 & 419 & 60.9 \\
S3 & sensing & determinants-late & off & 4888 & 11367 & 538 & 0 & 0 & 82.7 \\
S3 & sensing & determinants-late & on & 4393 & 11367 & 43 & 0 & 0 & 82.3 \\
\bottomrule
\end{tabular}}
\end{table}

\begin{table}[!htbp]\centering
\caption{Design variables and domain sizes of the four design-space
scales.}\label{tabS8}
{\let\scriptsize\tiny\setlength{\tabcolsep}{2pt}\scriptsize
\begin{tabular}{llrrrr}
\toprule
variable & meaning & S1 & S2 & S3 & S4 \\
\midrule
$\mathit{n\_e}$ & number of edge nodes & 3 & 6 & 24 & 80 \\
$\mathit{edge\_cls}$ & edge node class & 3 & 3 & 3 & 3 \\
$\mathit{n\_g}$ & number of gateways & 2 & 3 & 5 & 8 \\
$\mathit{gw\_cls}$ & gateway class & 2 & 2 & 2 & 2 \\
$\mathit{rep}$ & replication factor of stage 1 & 2 & 3 & 4 & 5 \\
$\mathit{policy}$ & placement policy & 3 & 3 & 3 & 3 \\
$\mathit{link\_cls}$ & access link class & 3 & 3 & 3 & 3 \\
$\mathit{cloud\_cls}$ & cloud instance class & 2 & 2 & 2 & 2 \\
\midrule
\multicolumn{2}{l}{\emph{design-space size} $|\mathcal{X}|$} & 1\,296 & 5\,832 & 51\,840 & 345\,600 \\
\bottomrule
\end{tabular}}
\end{table}

\begin{table}[!htbp]\centering
\caption{Workload classes. Stage values are given as
$\tau_1/\tau_2/\tau_3$.}\label{tabS9}
{\let\scriptsize\tiny\setlength{\tabcolsep}{2pt}\scriptsize
\begin{tabular}{lrrrrrrr}
\toprule
workload & $n_d$ & $\lambda$ & $\Lambda$ & $w_i$ [MI] & $s_i$ [kB] & $D$ [ms] & $C^{\max}$ \\
\midrule
telemetry & 200 & 2 & 400 & 8/4/20 & 20/10/10 & 120 & 60 \\
control & 60 & 10 & 600 & 5/3/9 & 5/4/3 & 30 & 60 \\
sensing & 40 & 1.5 & 60 & 120/30/260 & 1200/350/300 & 900 & 90 \\
\bottomrule
\end{tabular}}
\end{table}

\begin{table}[!htbp]\centering
\caption{Resource classes used in all experiments.}\label{tabS10}
{\let\scriptsize\tiny\setlength{\tabcolsep}{2pt}\scriptsize
\begin{tabular}{lrrrrr}
\toprule
class & $C$ / $B$ & $M$ / $\tau$ & $P^{\mathrm{idle}}$ / $e$ & $P^{\mathrm{peak}}$ & cost \\
\midrule
\multicolumn{6}{l}{\emph{compute node classes}} \\
edge-low & 2000 & 2048 & 4 & 12 & 1 \\
edge-med & 6000 & 8192 & 8 & 30 & 2.4 \\
edge-high & 15000 & 16384 & 15 & 65 & 5.2 \\
gw-low & 4000 & 4096 & 6 & 20 & 1.8 \\
gw-high & 12000 & 16384 & 12 & 48 & 4 \\
cloud-med & 40000 & 131072 & 0 & 180 & 6 \\
cloud-high & 100000 & 262144 & 0 & 380 & 12 \\
\midrule \multicolumn{6}{l}{\emph{link classes} (name, $B$ [Mb/s], $\tau$ [ms], $e$ [J/MB], cost)} \\
constrained & 10 & 4 & 0.3 & 0.2 & -- \\
standard & 100 & 2 & 0.12 & 0.6 & -- \\
fast & 1000 & 1 & 0.05 & 1.4 & -- \\
\bottomrule
\end{tabular}}
\end{table}

\begin{table}[!htbp]\centering
\caption{Exploration funnel across all scales and workloads.}\label{tabS11}
{\let\scriptsize\tiny\setlength{\tabcolsep}{2pt}\scriptsize
\setlength{\tabcolsep}{3.5pt}
\begin{tabular}{llrrrrrrrrrr}
\toprule
scale & workload & $|\mathcal{X}|$ & structural & feasibility & dominance & full & feasible & $|\mathcal{P}|$ & full-eval & exhaustive & recall \\
 & & & pruned & pruned & pruned & evals & & & reduction [\%] & evals & \\
\midrule
S1 & telemetry & 1\,296 & 216 & 208 & 840 & 32 & 25 & 17 & 97.53 & 1\,080 & 1.0 \\
S1 & control & 1\,296 & 216 & 226 & 826 & 28 & 25 & 18 & 97.84 & 1\,080 & 1.0 \\
S1 & sensing & 1\,296 & 216 & 530 & 534 & 16 & 11 & 11 & 98.77 & 1\,080 & 1.0 \\
S2 & telemetry & 5\,832 & 1\,860 & 965 & 2\,961 & 46 & 38 & 27 & 99.21 & 3\,888 & 1.0 \\
S2 & control & 5\,832 & 1\,860 & 1\,017 & 2\,916 & 39 & 36 & 25 & 99.33 & 3\,888 & 1.0 \\
S2 & sensing & 5\,832 & 1\,944 & 1\,534 & 2\,328 & 26 & 19 & 16 & 99.55 & 3\,888 & 1.0 \\
S3 & telemetry & 51\,840 & 12\,240 & 21\,957 & 17\,591 & 52 & 44 & 31 & 99.90 & 34\,560 & 1.0 \\
S3 & control & 51\,840 & 12\,240 & 21\,924 & 17\,633 & 43 & 40 & 27 & 99.92 & 34\,560 & 1.0 \\
S3 & sensing & 51\,840 & 14\,634 & 16\,230 & 20\,933 & 43 & 36 & 27 & 99.92 & 34\,560 & 1.0 \\
S4 & telemetry & 345\,600 & 33\,924 & 276\,742 & 34\,877 & 57 & 49 & 32 & 99.98 & -- & -- \\
S4 & control & 345\,600 & 33\,924 & 276\,811 & 34\,812 & 53 & 50 & 29 & 99.98 & -- & -- \\
S4 & sensing & 345\,600 & 41\,250 & 230\,523 & 73\,775 & 52 & 45 & 28 & 99.98 & -- & -- \\
\bottomrule
\end{tabular}}
\end{table}

\begin{table}[!htbp]\centering
\caption{Scalability sweep: expanded partial designs, full evaluations and
runtime.}\label{tabS12}
{\let\scriptsize\tiny\setlength{\tabcolsep}{2pt}\scriptsize
\setlength{\tabcolsep}{3.5pt}
\begin{tabular}{lrrrrrrrrr}
\toprule
& \multicolumn{3}{c}{telemetry} & \multicolumn{3}{c}{control} & \multicolumn{3}{c}{sensing} \\
\cmidrule(lr){2-4}\cmidrule(lr){5-7}\cmidrule(lr){8-10}
$|\mathcal{X}|$ & states & evals & [s] & states & evals & [s] & states & evals & [s] \\
\midrule
${1.3\times 10^{3}}$ & 286 & 42 & 0.05 & 268 & 34 & 0.05 & 234 & 18 & 0.04 \\
${5.8\times 10^{3}}$ & 696 & 54 & 0.14 & 682 & 54 & 0.13 & 649 & 30 & 0.10 \\
${5.2\times 10^{4}}$ & 2\,746 & 57 & 0.60 & 2\,732 & 57 & 0.52 & 3\,495 & 53 & 0.65 \\
${3.5\times 10^{5}}$ & 5\,092 & 58 & 1.14 & 5\,074 & 58 & 1.05 & 8\,228 & 59 & 1.63 \\
${2.8\times 10^{6}}$ & 8\,025 & 58 & 1.88 & 8\,007 & 58 & 1.70 & 13\,986 & 61 & 3.01 \\
${1.6\times 10^{7}}$ & 9\,252 & 58 & 2.22 & 9\,234 & 58 & 2.04 & 16\,345 & 61 & 3.67 \\
${8.3\times 10^{7}}$ & 10\,333 & 58 & 2.13 & 10\,315 & 58 & 2.50 & 18\,082 & 61 & 3.96 \\
\bottomrule
\end{tabular}}
\end{table}

\begin{table}[!htbp]\centering
\caption{Analytical mean latency versus the discrete-event evaluator: 40 designs per
workload, five seeds per design.}\label{tabS13}
{\let\scriptsize\tiny\setlength{\tabcolsep}{2pt}\scriptsize
\begin{tabular}{lrrr}
\toprule
workload & designs & MAPE [\%] & Spearman $\rho$ \\
\midrule
telemetry & 40 & 6.9 & 0.986 \\
control & 40 & 6.1 & 0.988 \\
sensing & 40 & 20.5 & 0.955 \\
\bottomrule
\end{tabular}}
\end{table}

\end{document}
